\documentclass[12pt,a4paper]{article}

% --------------------------------------------------
% Packages
% --------------------------------------------------
\usepackage[a4paper,margin=1.7cm]{geometry}
\usepackage{times}
\usepackage{array}
\usepackage{enumitem}
\usepackage{fancyhdr}
\usepackage{titlesec}
\usepackage{longtable}

% --------------------------------------------------
% Page Header / Footer
% --------------------------------------------------
\pagestyle{fancy}
\fancyhf{}

\fancyhead[L]{\textbf{Journal 3}}
\fancyhead[R]{BlurGuard -- Real-Time Screen-Share Sensitive Data Censor}

\renewcommand{\headrulewidth}{0.4pt}
\renewcommand{\footrulewidth}{0pt}

\fancyfoot[C]{\thepage}

% --------------------------------------------------
% Section Formatting
% --------------------------------------------------
\titleformat{\section}
  {\normalfont\bfseries\fontsize{11.5}{14}\selectfont}
  {\thesection.}{0.5em}{}

\titleformat{\subsection}
  {\normalfont\bfseries\fontsize{10}{12}\selectfont}
  {\thesubsection}{0.5em}{}

\titlespacing*{\section}
  {0pt}{7pt}{5pt}

\titlespacing*{\subsection}
  {0pt}{5pt}{3pt}

% --------------------------------------------------
% Paragraph / List Formatting
% --------------------------------------------------
\setlength{\parindent}{0pt}
\setlength{\parskip}{3pt}

\setlist[itemize]{
    leftmargin=1.2em,
    itemsep=1.5pt,
    topsep=2pt,
    parsep=0pt
}

% ==================================================
% DOCUMENT
% ==================================================
\begin{document}

% ==================================================
% JOURNAL HEADER
% ==================================================

\begin{center}

{\bfseries\fontsize{14}{16}\selectfont JOURNAL 3}

\vspace{3pt}

{\bfseries\fontsize{11}{13}\selectfont Frontend Implementation}

\vspace{3pt}

{\bfseries\fontsize{10.5}{13}\selectfont
BlurGuard -- Real-Time Screen-Share Sensitive Data Censor}

\end{center}

\vspace{5pt}

% ==================================================
% PROJECT INFORMATION TABLE
% ==================================================

\begin{center}
\renewcommand{\arraystretch}{1.15}

\begin{tabular}{|p{3.8cm}|p{12.4cm}|}
\hline

\textbf{Project}
&
BlurGuard -- Real-Time Screen-Share Sensitive Data Censor
\\
\hline

\textbf{Development Stage}
&
Frontend Prototype Implementation
\\
\hline

\textbf{Technology Used}
&
React, Vite, React Router, JavaScript, CSS
\\
\hline

\textbf{State Management}
&
React Context API
\\
\hline

\textbf{Student}
&
Lakshita
\\
\hline

\end{tabular}
\end{center}

\vspace{3pt}

% ==================================================
% 1. OBJECTIVE
% ==================================================

\section{Objective}

The objective of this development phase was to implement the frontend prototype of the
BlurGuard -- Real-Time Screen-Share Sensitive Data Censor system using React. The frontend
was developed to provide an interactive interface for users to configure privacy settings,
select the screen or application to monitor, and visualize real-time detection and protection
of sensitive information during screen sharing.

The work focused on application setup, navigation, user interaction flows, real-time
monitoring interface, configuration settings, and frontend integration of core features
to demonstrate the major workflow before backend and model integration.

% ==================================================
% 2. DEVELOPMENT ENVIRONMENT
% ==================================================

\section{Development Environment and Technologies}

The frontend was developed using the following technologies:

\begin{itemize}

    \item \textbf{React} -- used for developing reusable user interface components.

    \item \textbf{Vite} -- used for creating and running the React development project.

    \item \textbf{JavaScript} -- used for frontend application logic.

    \item \textbf{React Router} -- used for navigation between different application pages.

    \item \textbf{CSS} -- used for styling dashboards, forms, cards, navigation bars,
    buttons, and status indicators.

    \item \textbf{React Context API} -- used for maintaining shared application state
    between frontend components.

\end{itemize}

% ==================================================
% IMPORTANT:
% NO \newpage OR \clearpage HERE
% SECTION 3 WILL CONTINUE IN THE REMAINING SPACE
% ==================================================

% ==================================================
% 3. FRONTEND PROJECT SETUP
% ==================================================

\section{Frontend Project Setup}

The frontend project was organized as a modular React application. Separate components
and page sections were considered for authentication, dashboard monitoring, privacy settings,
detection results, and session controls. This structure allows individual parts of the
interface to be developed and modified without affecting the complete application.

The major frontend sections include:

\begin{itemize}

    \item Login / User Session

    \item Main Dashboard

    \item Screen and Application Selection

    \item Privacy Rules and Settings

    \item Live Monitoring

    \item Detection and Risk Results

    \item Protected Region / Censor Controls

    \item Alerts and User Review

\end{itemize}

% ==================================================
% 4. APPLICATION NAVIGATION
% ==================================================

\section{Application Navigation}

Frontend navigation was structured so that the user can move between the major BlurGuard
functions from a single application interface. The planned major routes and views are:

\begin{itemize}

    \item Login

    \item Dashboard

    \item Screen / Window Selection

    \item Privacy Settings

    \item Monitoring Session

    \item Detection Results

    \item Alerts / Review

    \item Session History / Logs

\end{itemize}

The navigation structure keeps the frontend organized and provides a clear path from
starting a session to monitoring and protecting sensitive information.

% ==================================================
% 5. LOGIN AND USER SESSION
% ==================================================

\section{Login and User Session Interface}

A simple login interface was included as the entry point of the BlurGuard frontend.
The interface is intended to establish the user session before the monitoring dashboard
is opened.

The interface contains:

\begin{itemize}

    \item Email / username field

    \item Password field

    \item Login button

    \item Session entry and logout controls

\end{itemize}

At the prototype stage, the interface focuses on frontend interaction. Actual authentication,
authorization, and secure session handling can be integrated with the backend in a later stage.

% ==================================================
% 6. DASHBOARD
% ==================================================

\section{BlurGuard Dashboard}

The main dashboard was designed as the central control interface for a screen-protection
session. It provides the user with a quick overview of the current monitoring state and
the actions required to start or manage protection.

The dashboard includes:

\begin{itemize}

    \item Monitoring status

    \item Selected screen or application

    \item Protection status

    \item Number of detected sensitive regions

    \item Current risk level

    \item Start / Pause / Stop monitoring controls

    \item Access to privacy settings

    \item Recent alerts or detection events

\end{itemize}

% ==================================================
% 7. SCREEN / APPLICATION SELECTION
% ==================================================

\section{Screen / Application Selection}

A dedicated interface was designed for selecting the screen, application window, or
display region that BlurGuard should monitor. This prevents unnecessary monitoring
of unrelated content and provides the user with control over the monitoring scope.

The selection interface can display available screens or windows and allow the user
to confirm the target before starting a monitoring session.

% ==================================================
% 8. PRIVACY RULES
% ==================================================

\section{Privacy Rules and Settings}

A privacy settings interface was designed to allow the user to configure which categories
of information should be protected. The interface can provide controls for sensitive
data categories and protection behaviour.

Example configuration options include:

\begin{itemize}

    \item Email addresses

    \item Phone numbers

    \item Passwords and authentication secrets

    \item API keys and access tokens

    \item Payment or card-related information

    \item Personal identification information

    \item QR codes or other sensitive visual content

    \item Blur, pixelate, or redact as the protection method

    \item Detection confidence threshold

\end{itemize}

% ==================================================
% 9. REAL-TIME MONITORING
% ==================================================

\section{Real-Time Monitoring Interface}

The monitoring interface was designed to represent the live state of the BlurGuard
system while a screen-sharing session is active. The interface provides visual feedback
so that the user can understand whether monitoring and protection are running.

The monitoring view can display:

\begin{itemize}

    \item Current monitoring state

    \item Selected screen or application

    \item Frames being monitored

    \item Detection activity

    \item Protected regions

    \item Current risk level

    \item Alerts requiring user attention

\end{itemize}

% ==================================================
% 10. DETECTION RESULTS
% ==================================================

\section{Detection Results and Alert Interface}

A detection-results section was designed to show information returned by the detection
and analysis layer. The purpose of this interface is to make the protection process
understandable to the user without exposing the sensitive information itself unnecessarily.

The interface can represent:

\begin{itemize}

    \item Detection category

    \item Confidence level

    \item Risk level

    \item Detected region

    \item Protection action applied

    \item Time of detection

\end{itemize}

An alert or review action can be provided when the system detects content that requires
confirmation or when the confidence level is below the configured threshold.

% ==================================================
% 11. FRONTEND COMPONENT STRUCTURE
% ==================================================

\section{Frontend Component Structure}

The frontend was organized around reusable components so that common interface
elements can be shared across different pages. A modular component structure also
makes future integration with backend APIs and real-time detection services easier.

\vspace{2pt}

\begin{center}
\renewcommand{\arraystretch}{1.15}

\begin{tabular}{|p{5.8cm}|p{10.4cm}|}
\hline

\textbf{Component / Page} & \textbf{Purpose}
\\
\hline

Login & User session entry
\\
\hline

Dashboard & Central monitoring controls and status
\\
\hline

ScreenSelector & Select screen or application to monitor
\\
\hline

PrivacySettings & Configure sensitive-data rules
\\
\hline

MonitoringPanel & Display live monitoring state
\\
\hline

DetectionPanel & Display detection and risk information
\\
\hline

AlertPanel & Display alerts and review actions
\\
\hline

ProtectionStatus & Show current censoring/protection state
\\
\hline

\end{tabular}
\end{center}

% ==================================================
% 12. STATE AND INTEGRATION
% ==================================================

\section{Frontend State and Integration}

Shared frontend state was considered for maintaining information such as the active
monitoring session, selected screen or application, privacy settings, detection results,
and protection status. React Context API can be used to make this information available
to multiple components without maintaining separate copies of the same state.

The frontend is structured so that the detection and analysis engine can later provide
real-time results through an API, WebSocket connection, or another application-level
communication mechanism. Similarly, the protected output can later be connected to
the screen-sharing integration layer.

% ==================================================
% 13. FRONTEND INTEGRATION AND TESTING
% ==================================================

\section{Frontend Integration and Testing}

The frontend prototype was checked from the perspective of navigation, interface
interaction, form behaviour, and consistency of the monitoring workflow.

The main testing activities include:

\begin{itemize}

    \item Starting the frontend application.

    \item Navigation between the main pages.

    \item Opening the monitoring dashboard.

    \item Selecting a screen or application.

    \item Changing privacy settings.

    \item Starting and stopping the monitoring interface.

    \item Displaying detection and risk information.

    \item Displaying protection status and alerts.

\end{itemize}

These checks help ensure that the frontend provides a consistent user flow before
connecting the interface to the complete detection and screen-sharing pipeline.

% ==================================================
% 14. CURRENT FRONTEND RESULT
% ==================================================

\section{Current Frontend Result}

The frontend stage establishes the user-facing structure of BlurGuard. The main dashboard,
screen/application selection, privacy configuration, monitoring controls, detection results,
and alert interfaces have been defined as the primary frontend workflow.

The current frontend prototype provides the foundation for integrating the actual OCR,
PII/secret detection, computer-vision analysis, protection engine, and screen-sharing
output in subsequent development stages.

% ==================================================
% 15. CURRENT LIMITATIONS
% ==================================================

\section{Current Limitations}

The current frontend stage has the following limitations:

\begin{itemize}

    \item Backend API integration is not yet part of this development stage.

    \item Actual authentication and authorization require backend support.

    \item Real-time detection results depend on integration with the detection engine.

    \item Actual screen capture and protected output require system-level integration.

    \item Persistent storage and complete session history require backend/database integration.

\end{itemize}

% ==================================================
% 16. LEARNING OUTCOMES
% ==================================================

\section{Learning Outcomes}

The frontend implementation provided practical understanding of:

\begin{itemize}

    \item React component-based development.

    \item Frontend routing and multi-page navigation.

    \item Designing interfaces for real-time monitoring systems.

    \item Interactive forms and configuration controls.

    \item Shared state management using React Context API.

    \item Modular frontend architecture.

    \item Frontend integration planning for backend and real-time services.

    \item Debugging and testing an interactive prototype before full system integration.

\end{itemize}

% ==================================================
% 17. RESULT
% ==================================================

\section{Result}

The following development activities were completed for this journal:

\vspace{2pt}

\begin{center}
\renewcommand{\arraystretch}{1.2}

\begin{tabular}{|p{10.8cm}|p{5.4cm}|}
\hline

\textbf{Activity} & \textbf{Status}
\\
\hline

Frontend project structure defined & Completed
\\
\hline

Main dashboard interface designed & Completed
\\
\hline

Screen / application selection interface & Completed
\\
\hline

Privacy settings interface & Completed
\\
\hline

Monitoring interface & Completed
\\
\hline

Detection and alert interface & Completed
\\
\hline

Frontend component structure & Completed
\\
\hline

Frontend workflow testing & Completed
\\
\hline

\end{tabular}
\end{center}

% ==================================================
% 18. CONCLUSION
% ==================================================

\section{Conclusion}

Journal 03 documents the frontend implementation stage of the BlurGuard --
Real-Time Screen-Share Sensitive Data Censor project. The frontend structure
establishes the main user-facing workflow for starting a monitoring session,
selecting a screen or application, configuring privacy rules, viewing monitoring
and detection information, and controlling the protection process.

The completed frontend design provides a foundation for the next development
stages, including integration with the OCR and sensitive-data detection engine,
context and risk analysis, sensitive-region tracking, protection engine, and
screen-sharing output.

\end{document}